\section{Context (\texttt{env.context})}

\begin{frame}[fragile]{What Is the Context?}
	\begin{itemize}
		\item \textbf{\texttt{env.context}} is a dictionary of flags and values for the current operation.
\begin{lstlisting}
print(self.env.context)
# example: {'lang': 'en_US', 'tz': 'Africa/Mogadishu', 'uid': 2}
\end{lstlisting}
		\item Context can change:
		\begin{itemize}
			\item language
			\item default values
			\item active test behavior
			\item custom business flags
		\end{itemize}
	\end{itemize}
\end{frame}

\begin{frame}[fragile]{Reading Context Values}
\begin{lstlisting}
lang = self.env.context.get('lang')
force_email = self.env.context.get('force_send_email')
default_gender = self.env.context.get('default_gender')
\end{lstlisting}
	\begin{itemize}
		\item Use \texttt{.get()} so missing keys do not crash your code.
		\item Keys starting with \texttt{default\_} are used by \texttt{default\_get()}.
	\end{itemize}
\end{frame}

\begin{frame}[fragile]{\texttt{with\_context()} Method}
\begin{block}{Method Syntax}
\begin{lstlisting}
records.with_context(key=value, ...)
# or
records.with_context({...})
\end{lstlisting}
\end{block}
\begin{lstlisting}
students = self.env['student.student'].with_context(
	default_gender='female',
	active_test=False,
)
\end{lstlisting}
	\begin{itemize}
		\item Returns a new recordset/environment with updated context.
		\item Original environment is not modified.
	\end{itemize}
\end{frame}

\begin{frame}[fragile]{Common Context Keys}
	\begin{itemize}
		\item \texttt{lang} -- language code
		\item \texttt{tz} -- timezone
		\item \texttt{active\_test} -- whether to hide inactive records
		\item \texttt{default\_fieldname} -- default value for create forms
		\item custom keys for your own business logic
	\end{itemize}
\begin{lstlisting}
self.with_context(skip_student_validation=True).write(vals)
\end{lstlisting}
\end{frame}
